\documentclass[11pt,a4paper]{article}

\usepackage[margin=1in]{geometry}
\usepackage{amsmath, amssymb, amsthm}
\usepackage{mathtools}
\usepackage{bm}
\usepackage{enumitem}
\usepackage{microtype}
\usepackage[T1]{fontenc}
\usepackage{lmodern}
\usepackage[dvipsnames]{xcolor}
\usepackage{hyperref}

\hypersetup{
  colorlinks=true,
  linkcolor=blue!50!black,
  urlcolor=blue!50!black,
  citecolor=blue!50!black
}

\newtheorem{theorem}{Theorem}[section]
\newtheorem{lemma}[theorem]{Lemma}
\newtheorem{definition}[theorem]{Definition}

\title{U3: The Forced Structure of Time-Translation Symmetry and Schr\"odinger Evolution}

\author{Allen Proxmire}

\date{April 2026}

\begin{document}

\maketitle

\begin{abstract}
We present the structural derivation of the participation-measure evolution equation $i\hbar \, \partial_t P_K = \hat{H} P_K$ --- the Schr\"odinger equation in standard non-relativistic single-particle form --- from Primitive~13 (relational timing), Primitive~03 (graph homogeneity), and the now-\textsc{forced} upstream items U1, U2-Discrete, U2-Continuum, U5, and U4. Operating under a strict identification-not-derivation discipline (U4 may appear in the F5 step only as the identification target, never as a derivation premise), we decompose U3 into five sub-features: existence of a strongly continuous time-translation representation on the U2 Hilbert space (F1); uniqueness of its self-adjoint generator via Stone's theorem (F2); linearity of the induced evolution equation (F3); first-order-in-time structure (F4); and compatibility of the Stone-derived generator $\hat{H}_{U3}$ with the U4-derived kinetic-plus-potential operator $\hat{H}_{U4} = \hbar^2|\hat{p}|^2/(2m) + V(\hat{x})$ (F5). F1 follows from Primitive~13 (continuous time axis) plus Primitive~03 (homogeneity) plus U2 inner-product translation-invariance. F2 follows from Stone's theorem on one-parameter unitary groups. F3 and F4 are automatic from Hilbert-space structure and the operator-exponential parameterization. F5 --- the load-bearing content --- is established by \textbf{Galilean Lie algebra closure}: $\hat{H}_{U3}$, as the time-translation generator within the Galilean representation on the U2 Hilbert space, must satisfy the commutator $[\hat{H}_{U3}, \hat{K}_i] = -i\hbar\hat{p}_i$ with the boost generator $\hat{K}_i = m\hat{x}_i - t\hat{p}_i$; integrating against the U5 momentum operator yields $T(\hat{p}) = |\hat{p}|^2/(2m)$ via the chain-rule Jacobian, identifying $\hat{H}_{U3}$ with $\hat{H}_{U4}$ modulo the additive vacuum-energy zero-point gauge. The identification is cross-checked by Stone-theorem uniqueness applied to the $\hat{H}_{U4}$-generated unitary group. All seven falsifiers are dispatched. The U3 verdict is \textbf{\textsc{forced}}, with no new structural \textsc{candidate}s introduced and zero new conditionalities. \textbf{U4 is retroactively promoted from \textsc{forced}-conditional to \textsc{forced}-unconditional}, resolving its Framing-1 sibling-\textsc{candidate} conditionality. The active upstream-\textsc{candidate} inventory of the QM-emergence Phase-1 program reduces from $\{U3\}$ + continuum gauge to $\{ \, \}$ + continuum gauge --- empty modulo description-level convention. \textbf{All four foundational postulates of non-relativistic single-particle quantum mechanics --- the Born rule, the Bell--Tsirelson bound, the Heisenberg uncertainty inequality, and the Schr\"odinger evolution equation --- are now structurally \textsc{forced}-unconditional consequences of the Event Density primitive ontology.} The forced-theorem inventory rises to sixteen with two retroactive promotions; Phase-1 of the QM-emergence program is structurally complete.
\end{abstract}

\section{Introduction}

\subsection{What U3 is}

U3 is the structural commitment that the participation-measure components $P_K$ on the U2 Hilbert space $\mathcal{H}$ evolve in time according to a single-generator linear first-order equation,
\[
i\hbar \, \partial_t P_K(x, t) = \hat{H} P_K(x, t),
\]
with $\hat{H}$ a self-adjoint operator on $\mathcal{H}$. Combined with U4's specific Hamiltonian form $\hat{H} = \hbar^2|\hat{p}|^2/(2m) + V(\hat{x})$, the equation is the standard Schr\"odinger equation of non-relativistic single-particle quantum mechanics. U3 supplies the existence, uniqueness, linearity, and first-order character of the equation; U4 supplies the form of the generator $\hat{H}$.

The earlier framing of U3 in the QM-emergence Phase-1 inventory bundled it with U4 as a single dynamical commitment. The present arc separates the two cleanly: U4 was closed first (in the 2026-04-26 cycle, conditional on U3 via Framing~1); U3 is now closed independently (this paper), retroactively resolving U4's conditionality. The two-arc separation is methodologically essential to preserve acyclicity in the U3 $\leftrightarrow$ U4 loop.

\subsection{Why U3 matters}

U3 is \textbf{the final arc of the QM-emergence Phase-1 structural-foundations program}. With U3 closed, every upstream commitment in the participation-measure $\to$ Hilbert-space $\to$ Schr\"odinger-emergence chain has been derived from the ED primitive ontology plus standard mathematical and symmetry-group infrastructure. The four foundational postulates of non-relativistic single-particle quantum mechanics --- the Born rule (Theorem~10), the Bell--Tsirelson bound, the Heisenberg uncertainty inequality, and the Schr\"odinger evolution equation --- are now \textbf{\textsc{forced}-unconditional structural consequences of ED primitives}.

The result is the closure of a seven-arc cycle (born\_gleason $\to$ U2-Discrete $\to$ U2-Continuum $\to$ U5 $\to$ U1 $\to$ U4 $\to$ U3) in which the structural-derivation methodology has been demonstrated against seven substantively different structural questions. Phase-1 is structurally complete, modulo the description-level continuum gauge convention established in U2-Continuum.

\subsection{Relationship to U1, U2, U5, U4}

\begin{itemize}
  \item \textbf{U1.} U3 takes the participation-measure construction $P_K = \sqrt{b_K} \cdot e^{i\pi_K}$ as input. The $U(1)$-valued polarity-phase factor supplies the complex structure on which Hermitian operators act and on which time evolution preserves the inner product.
  \item \textbf{U2.} U3 takes the U2 Hilbert space $\mathcal{H}$ and its sesquilinear inner product as input. Time-translation operators are unitary on $\mathcal{H}$; Stone's theorem applies.
  \item \textbf{U5.} U3 takes the U5 momentum operator $\hat{p} = -i\hbar\nabla$ as input for the F5 Galilean-closure step. The U5 derivation also serves as the methodological template for the U3 arc: Stone's theorem applied to time-translation symmetry mirrors the U5 application to spatial-translation symmetry.
  \item \textbf{U4.} U3 treats U4's operator form $\hat{H}_{U4} = \hbar^2|\hat{p}|^2/(2m) + V(\hat{x})$ as the \textbf{identification target} in F5, never as a derivation premise. The strict discipline that distinguishes ``U4 as identification target'' from ``U4 as derivation input'' preserves acyclicity; without this discipline, U3's derivation would invert U4's Framing-1 conditionality and close a circular loop.
\end{itemize}

\subsection{Summary of the U3 verdict}

The U3 verdict is \textbf{\textsc{forced}}. The five sub-features F1, F2, F3, F4, F5 close uniformly under the available structural inputs and the identification-not-derivation discipline. All seven falsifiers (Fal-1 through Fal-7) dispatch. No new structural \textsc{candidate}s are introduced. \textbf{U4 is retroactively promoted to \textsc{forced}-unconditional}, completing the structural derivation of the Schr\"odinger equation.

The downstream cascade is decisive:

\begin{itemize}
  \item The active upstream-\textsc{candidate} inventory reduces from $\{U3\}$ + continuum gauge to $\{ \, \}$ + continuum gauge.
  \item The forced-theorem count rises from fifteen to sixteen, with two retroactive promotions in this cycle (U3 from \textsc{candidate} to \textsc{forced}; U4 from \textsc{forced}-conditional to \textsc{forced}-unconditional).
  \item All four foundational QM postulates are \textsc{forced}-unconditional.
  \item \textbf{Phase-1 of the QM-emergence program is structurally complete} (modulo description-level continuum gauge).
\end{itemize}

\subsection{Organization}

Section~\ref{sec:question} states the U3 question precisely and inventories its sub-features and falsifiers. Section~\ref{sec:inputs} inventories the structural inputs and the inputs explicitly forbidden under the identification-not-derivation discipline. Section~\ref{sec:decomposition} decomposes U3 into F1--F5 with status assignments. Section~\ref{sec:falsifiers} introduces the seven falsifiers. Sections~\ref{sec:F1}--\ref{sec:F4} derive F1, F2, F3, F4 in turn. Section~\ref{sec:F5} derives the load-bearing F5 via Galilean Lie algebra closure (Strategy A) with Stone-uniqueness as cross-check (Strategy B). Section~\ref{sec:falsifier-table} consolidates the falsifier dispatches. Section~\ref{sec:conditionality} documents the conditionality ledger. Section~\ref{sec:final-theorem} states the final U3 theorem. Section~\ref{sec:downstream} documents the downstream cascade including Phase-1 completion. Section~\ref{sec:discussion} discusses the result's significance and identifies the next-priority work.

\section{Statement of the U3 Question}\label{sec:question}

\subsection{The target structure}

The target structural content of U3 is the existence, uniqueness, linearity, and first-order-in-time character of the participation-measure evolution equation, together with the identification of its generator with the U4-derived Hamiltonian operator.

Formally: let $\mathcal{H}$ denote the U2 Hilbert space of participation measures (discrete or continuum regime, gauge-class equivalent under the U2-Continuum conformal redundancy). Let $U_t$ for $t \in \mathbb{R}$ denote the time-translation operator $(U_t P)_K(x, s) := P_K(x, s - t)$ acting componentwise on participation-measure components. The U3 question is whether $\{U_t : t \in \mathbb{R}\}$ is a strongly continuous one-parameter unitary group whose unique self-adjoint generator $\hat{H}_{U3}$ coincides with the U4-derived operator $\hat{H}_{U4} = \hbar^2|\hat{p}|^2/(2m) + V(\hat{x})$, yielding the Schr\"odinger evolution equation
\begin{equation}
i\hbar \, \partial_t P(t) = \hat{H} P(t), \qquad \hat{H} = \frac{\hbar^2|\hat{p}|^2}{2m} + V(\hat{x}). \label{eq:U3target}
\end{equation}

\subsection{The five sub-features}

\begin{itemize}
  \item \textbf{F1 (Existence of strongly continuous time-translation representation).} $\{U_t : t \in \mathbb{R}\}$ is a well-defined, strongly continuous, one-parameter abelian unitary group on $\mathcal{H}$.
  \item \textbf{F2 (Uniqueness of the self-adjoint generator).} There exists a unique self-adjoint operator $\hat{H}_{U3}$ on $\mathcal{H}$ such that $U_t = \exp(-i\hat{H}_{U3} t/\hbar)$.
  \item \textbf{F3 (Linearity of the evolution equation).} The induced equation $i\hbar \, \partial_t P = \hat{H}_{U3} P$ is linear in $P$.
  \item \textbf{F4 (First-order-in-time structure).} The evolution equation is first-order in $\partial_t$.
  \item \textbf{F5 (Compatibility with U4 Hamiltonian form).} The Stone-derived generator $\hat{H}_{U3}$ coincides as an operator with the U4-derived $\hat{H}_{U4} = \hbar^2|\hat{p}|^2/(2m) + V(\hat{x})$.
\end{itemize}

\subsection{The seven falsifiers}

\begin{itemize}
  \item \textbf{Fal-1 (Non-unitary time evolution).}
  \item \textbf{Fal-2 (Multiple inequivalent generators).}
  \item \textbf{Fal-3 (Non-linear evolution).}
  \item \textbf{Fal-4 (Higher-order-in-time evolution).}
  \item \textbf{Fal-5 (Generator incompatible with U4).}
  \item \textbf{Fal-6 (Non-self-adjoint generator).}
  \item \textbf{Fal-7 (Non-continuous time-translation representation).}
\end{itemize}

\subsection{Methodological framing: identification, not derivation}

The U3 arc adopts a strict discipline distinguishing two roles for U4:

\begin{itemize}
  \item \textbf{Forbidden:} U4 as a derivation premise. Using U4's operator form to derive $\hat{H}_{U3}$ would invert U4's Framing-1 conditionality (U4 was derived conditional on U3's existence-of-Hermitian-generator content) and close a circular loop.
  \item \textbf{Permitted:} U4 as the identification target in F5. The F5 step \emph{identifies} $\hat{H}_{U3}$ with $\hat{H}_{U4}$ via the shared Galilean Lie algebra structure --- both operators land on the same Hamiltonian because the Galilean group's representation theory forces them to, not because either is derived from the other.
\end{itemize}

This protocol preserves acyclicity in the U3 $\leftrightarrow$ U4 loop. The discipline is enforced by an explicit circularity audit (\S\ref{sec:forbidden}).

\section{Structural Inputs}\label{sec:inputs}

\subsection{Primitives}

\begin{itemize}
  \item \textbf{Primitive 13 (Relational Timing).} Supplies the continuous time axis $\mathbb{R}_t$ and time-translation symmetry as a kinematic feature of the participation graph. Central new input of this arc.
  \item \textbf{Primitive 03 (Participation).} Graph homogeneity, including time-axis homogeneity (no privileged instant of time).
  \item \textbf{Primitive 04 (Bandwidth).} Supplies the participation-measure components $P_K$ on which $U_t$ acts.
  \item \textbf{Primitive 09 (Tension Polarity).} $U(1)$-valued polarity phase. Supplies the complex structure on which Hermitian operators act (via U1).
\end{itemize}

\subsection{Inherited \textsc{forced} upstream items}

\begin{itemize}
  \item \textbf{U1 (Theorem 14).} Participation-measure construction $P_K = \sqrt{b_K} \cdot e^{i\pi_K}$.
  \item \textbf{U2-Discrete (Theorem 11) and U2-Continuum (Theorem 12).} Sesquilinear inner product on the participation-measure Hilbert space, in both regimes.
  \item \textbf{U5 (Theorem 13).} Spatial translation generator $\hat{p} = -i\hbar\nabla$. Methodological template; supplies the momentum operator entering F5.
  \item \textbf{U4 (Theorem 15).} Specific Hamiltonian form $\hat{H}_{U4} = \hbar^2|\hat{p}|^2/(2m) + V(\hat{x})$. \textbf{Used solely in F5 as the identification target.}
  \item \textbf{Theorem 10 (Born rule, Gleason--Busch).} Available as background.
\end{itemize}

\subsection{Mathematical infrastructure}

\begin{itemize}
  \item Stone's theorem on one-parameter unitary groups (Stone 1932)~\cite{stone}.
  \item Spectral theorem for self-adjoint operators on Hilbert spaces.
  \item Standard $L^2$ functional analysis: translation continuity, operator-exponential calculus~\cite{reedsimon}.
  \item Galilean Lie algebra and central extension (Bargmann 1954)~\cite{bargmann}.
\end{itemize}

\subsection{Forbidden inputs (circularity audit)}\label{sec:forbidden}

The following inputs are explicitly forbidden as derivation inputs to U3:

\begin{itemize}
  \item \textbf{U4 as derivation premise} for F1, F2, F3, F4. U4 may appear in F5 only as identification target, never as a derivation input.
  \item \textbf{U4's Framing-1 ``U3 exists'' assumption.} Forbidden as input to F5; F5 treats U4's operator form as a mathematical object on $\mathcal{H}$, not as a conditional construct.
  \item \textbf{U3 itself} (question-begging). Forbidden in any role.
  \item \textbf{The Schr\"odinger equation in textbook form.} Forbidden as derivation input.
  \item \textbf{Relativistic kinematics.} Out of scope.
  \item \textbf{Downstream gauge-coupling structure.} Out of scope (gauge-coupling-free per inherited Phase-1 Step~2 scope).
  \item \textbf{QFT-level inputs.} Out of scope.
\end{itemize}

Each subsequent derivation operates strictly within the primitive-input scope of the present arc. The circularity audit follows the canonical templates of U1 Memo~01 \S4, U4 Memo~01 \S2, and U5 \S3.5.

\subsection{Scope conditions}

\begin{itemize}
  \item \textbf{Non-relativistic single-particle scope.} Inherited from Phase-1 Step~2; selects the abelian $\mathbb{R}_t$ time-translation group and the Galilean Lie algebra.
  \item \textbf{Gauge-coupling-free scope.} Inherited from U4.
  \item \textbf{Flat-spacetime baseline (continuum regime).}
  \item \textbf{Continuous-time interpretation of Primitive~13.} Confirmed by pre-derivation audit at Phase-1 scope.
\end{itemize}

\section{Sub-Feature Decomposition}\label{sec:decomposition}

\begin{center}
\begin{tabular}{p{2.5cm}p{2.0cm}p{6.5cm}p{2.5cm}}
\hline
Sub-feature & Status & Structural inputs & Falsifiers \\
\hline
F1 (translation rep.) & \textsc{forced} & P-13 + P-03 + U2 inner-product translation-invariance & Fal-1, Fal-7 \\
F2 (unique generator) & \textsc{forced} & F1 + Stone's theorem & Fal-2, Fal-6 \\
F3 (linearity) & automatic & F1 + F2 + Hilbert-space operator structure & Fal-3 \\
F4 (first-order-in-$t$) & automatic & F2 + non-rel scope & Fal-4 \\
F5 (compat.\ with U4) & \textsc{forced} (load-bearing) & F1 + F2 + U5 momentum + Galilean Lie algebra; Stone-uniqueness cross-check & Fal-5 \\
\hline
\end{tabular}
\end{center}

Two forced-via-derivation items (F1, F2), two automatic items (F3, F4), one load-bearing item (F5).

\section{Falsifier Inventory}\label{sec:falsifiers}

\begin{itemize}
  \item \textbf{Fal-1 (Non-unitary time evolution).} $U_t$ fails to preserve the U2 inner product. Resolved \S\ref{sec:F1}.4.
  \item \textbf{Fal-2 (Multiple inequivalent generators).} Stone's theorem fails. Resolved \S\ref{sec:F2}.2.
  \item \textbf{Fal-3 (Non-linear evolution).} Resolved \S\ref{sec:F3}.2.
  \item \textbf{Fal-4 (Higher-order-in-time).} Klein--Gordon-like second-order equation. Resolved \S\ref{sec:F4}.2.
  \item \textbf{Fal-5 (Generator incompatible with U4).} $\hat{H}_{U3} \neq \hat{H}_{U4}$. Resolved \S\ref{sec:F5}.
  \item \textbf{Fal-6 (Non-self-adjoint generator).} Resolved \S\ref{sec:F2}.3.
  \item \textbf{Fal-7 (Non-continuous time-translation representation).} Resolved \S\ref{sec:F1}.6.
\end{itemize}

\section{Derivation of F1 (Time-Translation Representation)}\label{sec:F1}

\subsection{The claim}

The time-translation operators $\{U_t : t \in \mathbb{R}\}$ acting on $\mathcal{H}$ by
\[
(U_t P)_K(x, s) := P_K(x, s - t)
\]
form a strongly continuous one-parameter abelian unitary group.

\subsection{Time-translation symmetry as a kinematic feature}

Three primitive-level observations:

\begin{itemize}
  \item[\textbf{(i)}] Primitive~13 supplies a continuous time axis $\mathbb{R}_t$ at the structural level.
  \item[\textbf{(ii)}] Primitive~03's homogeneity extends to the time axis: the participation graph's structure at time $s$ is structurally identical to its structure at time $s - t$, up to relabeling.
  \item[\textbf{(iii)}] No primitive-level structure singles out a privileged time origin or magnitude.
\end{itemize}

These observations establish time-translation symmetry as a \emph{kinematic} property of the participation graph. \textbf{No reference to U4, no Hamiltonian, and no Schr\"odinger equation appears in this argument.}

\subsection{Linearity}

Time translation acts linearly on participation-measure components:
\[
U_t (\alpha P + \beta Q)_K(x, s) = \alpha (U_t P)_K(x, s) + \beta (U_t Q)_K(x, s)
\]
for all $\alpha, \beta \in \mathbb{C}$ and $P, Q \in \mathcal{H}$.

\subsection{Unitarity}

For $P, Q \in \mathcal{H}$ and $t \in \mathbb{R}$:
\[
\langle U_t P \mid U_t Q \rangle_{\mathcal{H}} = \sum_K \int_{\mathbb{R}^d} \int_{\mathbb{R}_t} P_K^*(x, s - t) \, Q_K(x, s - t) \, ds \, d\mu(x).
\]
Substituting $s' = s - t$, with $ds' = ds$ (the U2 measure is translation-invariant on the time axis: Primitive~13 supplies translation-invariant relational time, with no $t$-dependent weighting):
\[
\langle U_t P \mid U_t Q \rangle_{\mathcal{H}} = \langle P \mid Q \rangle_{\mathcal{H}}.
\]
\textbf{Falsifier Fal-1 dispatched.}

\subsection{Group structure}

Direct calculation: $U_t U_s = U_{t+s}$, $U_0 = I$, $U_{-t} = U_t^{-1}$. Abelian one-parameter group.

\subsection{Strong continuity}

For any $P \in \mathcal{H}$ with components $P_K(x, \cdot) \in L^2(\mathbb{R}_t, ds)$, the map $t \mapsto U_t P$ is strongly continuous because $L^2$-translation is strongly continuous (standard result of $L^2$ functional analysis~\cite{reedsimon}). Continuity inherits to the U2 Hilbert space's direct-integral structure.

The continuous-time interpretation of Primitive~13 is the precondition. Pre-derivation audit confirms continuous time at Phase-1 scope. \textbf{Falsifier Fal-7 dispatched.}

\begin{theorem}[F1 --- Time-Translation Representation]
Let $\mathcal{H}$ denote the U2 Hilbert space of participation measures. Then the time-translation operators $(U_t P)_K(x, s) := P_K(x, s - t)$ for $t \in \mathbb{R}$ form a strongly continuous one-parameter abelian unitary group on $\mathcal{H}$. The derivation invokes Primitives~13 and 03 plus the U2 inner-product structure, with no input from U4, no Hamiltonian, and no PDE-level Schr\"odinger assumption.
\end{theorem}

\section{Derivation of F2 (Stone-Uniqueness)}\label{sec:F2}

\subsection{Stone's theorem}

\begin{theorem}[Stone 1932~\cite{stone}]
Every strongly continuous one-parameter unitary group $\{V(t) : t \in \mathbb{R}\}$ on a Hilbert space has a unique self-adjoint generator $A$ such that $V(t) = \exp(itA)$.
\end{theorem}

\subsection{Application to F1's group}

By Theorem~F1, $\{U_t : t \in \mathbb{R}\}$ is a strongly continuous one-parameter abelian unitary group on $\mathcal{H}$. Apply Stone's theorem: there exists a unique self-adjoint operator $\hat{H}_{U3}$ on $\mathcal{H}$ such that
\begin{equation}
U_t = \exp\!\left( -\frac{i \hat{H}_{U3} t}{\hbar} \right). \label{eq:Ut-stone}
\end{equation}
The generator is unique. \textbf{Falsifier Fal-2 dispatched.}

\subsection{Self-adjointness}

Stone's theorem produces a self-adjoint generator from a strongly continuous unitary group. Self-adjointness is automatic. \textbf{Falsifier Fal-6 dispatched.}

\subsection{No U4-dependent assumptions enter}

The derivation invokes only F1 and Stone's theorem. U4's specific operator form does not enter. The operator $\hat{H}_{U3}$ is \emph{abstractly} the time-translation generator; identification with U4's operator is the F5 question.

\begin{theorem}[F2 --- Uniqueness of the Time-Translation Generator]
Under the hypotheses of Theorem~F1, there exists a unique self-adjoint operator $\hat{H}_{U3}$ on $\mathcal{H}$ such that $U_t = \exp(-i\hat{H}_{U3} t/\hbar)$ for all $t \in \mathbb{R}$. The identification is U4-independent.
\end{theorem}

\section{Derivation of F3 (Linearity)}\label{sec:F3}

\subsection{The argument}

Differentiating the operator-exponential at $t = 0$:
\[
\frac{d}{dt}\bigg|_{t=0} U_t P = -\frac{i}{\hbar} \hat{H}_{U3} P.
\]
For the time-evolved state $P(t) := U_t P_0$:
\begin{equation}
i\hbar \, \partial_t P(t) = \hat{H}_{U3} P(t). \label{eq:U3-eqn}
\end{equation}

\subsection{Linearity is automatic}

Equation~\eqref{eq:U3-eqn} is linear in $P(t)$ because: (i) $\mathcal{H}$ is a complex vector space; (ii) $\hat{H}_{U3}$ is a self-adjoint operator on $\mathcal{H}$, hence linear; (iii) $\{U_t\}$ acts linearly on $\mathcal{H}$ (per F1); (iv) time translations form an additive abelian group with $U_{t+s} = U_t U_s$ inheriting linearity.

Gross--Pitaevskii-like nonlinearities are phenomenological; they do not arise structurally in the Phase-1 single-particle scope. \textbf{Falsifier Fal-3 dispatched.}

\begin{theorem}[F3 --- Linearity]
The induced evolution equation $i\hbar \, \partial_t P(t) = \hat{H}_{U3} P(t)$ is linear in $P(t)$. Linearity is automatic from the Hilbert-space structure, the operator-theoretic definition of $\hat{H}_{U3}$, the linearity of the unitary representation $\{U_t\}$, and the additive structure of the time-translation group.
\end{theorem}

\section{Derivation of F4 (First-Order-in-Time)}\label{sec:F4}

\subsection{The argument}

Differentiating $U_t = \exp(-i\hat{H}_{U3} t/\hbar)$ once in $t$ produces a first-order equation \eqref{eq:U3-eqn}.

\subsection{Why higher-order does not arise}

Higher derivatives produce powers of $\hat{H}_{U3}$:
\[
i\hbar \, \partial_t^n P(t) = (\hat{H}_{U3})^n P(t),
\]
which are downstream consequences of \eqref{eq:U3-eqn}, not independent equations.

A second-order-in-time equation (Klein--Gordon-like) arises in \emph{relativistic} quantum mechanics. The non-relativistic scope condition (inherited from Phase-1 Step~2's framing, not from U4) selects the abelian $\mathbb{R}_t$ time-translation group; under this scope, the operator-exponential structure gives a first-order equation. \textbf{Falsifier Fal-4 dispatched.}

\begin{theorem}[F4 --- First-Order-in-Time]
The evolution equation $i\hbar \, \partial_t P(t) = \hat{H}_{U3} P(t)$ is first-order in $\partial_t$. The first-order character is automatic from Stone's theorem, the abelian group structure of time translations, and the non-relativistic scope condition.
\end{theorem}

\section{Derivation of F5 (Compatibility with U4)}\label{sec:F5}

This is the load-bearing derivation. It establishes $\hat{H}_{U3} = \hat{H}_{U4}$ via Galilean Lie algebra closure (Strategy A), with Stone-theorem uniqueness (Strategy B) as cross-check.

\subsection{The exact compatibility condition}

After F1--F4, two self-adjoint operators are in play on $\mathcal{H}$:

\begin{itemize}
  \item $\hat{H}_{U3}$: the unique self-adjoint generator of $\{U_t\}$, abstractly the ``time-translation generator.''
  \item $\hat{H}_{U4} = \hbar^2|\hat{p}|^2/(2m) + V(\hat{x})$: derived in U4 via Galilean integration on the U5 momentum operator.
\end{itemize}

The F5 question:
\begin{equation}
\hat{H}_{U3} \stackrel{?}{=} \hat{H}_{U4} \quad \text{as self-adjoint operators on } \mathcal{H}. \label{eq:F5-question}
\end{equation}

\subsection{Discipline: identification, not derivation}

Per \S\ref{sec:forbidden}, the F5 question is not ``derive $\hat{H}_{U3}$ from $\hat{H}_{U4}$'' (which would close the U3 $\leftrightarrow$ U4 circular loop). The question is whether the two operators, each constructed independently, are the same operator on $\mathcal{H}$. The protocol:

\begin{itemize}
  \item $\hat{H}_{U3}$ is derived from F1 + F2 (Stone's theorem on time translations). U4 plays no role.
  \item $\hat{H}_{U4}$ is treated as a mathematical object on $\mathcal{H}$. U4's derivation chain is taken as background.
  \item F5 establishes an identification via shared structural infrastructure (the Galilean Lie algebra), not via deriving one from the other.
\end{itemize}

\subsection{The Galilean Lie algebra at primitive level}

The Galilean Lie algebra is generated by $\hat{H}$ (time translations), $\hat{p}_i$ (spatial translations; the U5 momentum operator), $\hat{x}_i$ (position), $\hat{K}_i$ (Galilean boosts), and $\hat{J}_i$ (rotations). The substantive commutators are:
\begin{align}
[\hat{x}_i, \hat{p}_j] &= i\hbar \, \delta_{ij}, \label{eq:xp-comm} \\
[\hat{K}_i, \hat{p}_j] &= i\hbar m \, \delta_{ij}, \label{eq:Kp-comm} \\
[\hat{H}, \hat{K}_i] &= -i\hbar \, \hat{p}_i. \label{eq:HK-comm}
\end{align}

The mass $m$ enters \eqref{eq:Kp-comm} as the Galilean group's central-extension structure constant (Bargmann 1954)~\cite{bargmann}. The boost generator takes the form
\begin{equation}
\hat{K}_i = m \, \hat{x}_i - t \, \hat{p}_i, \label{eq:K-form}
\end{equation}
forced by the requirement that $\hat{K}_i$ generate the standard Galilean boost $\hat{p} \to \hat{p} - mv$.

The Galilean Lie algebra is \emph{primitive-level non-relativistic-scope infrastructure}, not specific to U4's derivation. U4 \emph{uses} the algebra to derive $\hat{H}_{U4}$'s form; U3 \emph{uses} the same algebra (independently) to identify $\hat{H}_{U3}$ as the time-translation generator within the algebra's representation.

\subsection{The U2 Galilean representation}

All Galilean transformations act unitarily on $\mathcal{H}$:

\begin{itemize}
  \item Spatial translations $T_a = \exp(i\hat{p}\cdot a/\hbar)$: unitary by U5's F1.
  \item Time translations $U_t = \exp(-i\hat{H}_{U3} t/\hbar)$: unitary by U3's F1.
  \item Rotations $R_\omega$: unitary by primitive-level rotation invariance.
  \item Galilean boosts $B_v = \exp(i\hat{K}\cdot v/\hbar)$: unitary as exponential of self-adjoint $\hat{K}$.
\end{itemize}

The U2 Hilbert space carries a unitary representation of the Galilean group.

\subsection{Strategy A --- Galilean Lie algebra closure}

By F1 + F2, $\hat{H}_{U3}$ is the unique self-adjoint generator of $\{U_t\}$ on $\mathcal{H}$. By \S\ref{sec:F5}.4, $\{U_t\}$ is part of the Galilean group's unitary representation on $\mathcal{H}$. Therefore $\hat{H}_{U3}$ is \textbf{the time-translation generator within the Galilean representation} --- it plays the role of $\hat{H}$ in \eqref{eq:xp-comm}--\eqref{eq:HK-comm}. In particular, $\hat{H}_{U3}$ must satisfy \eqref{eq:HK-comm}:
\begin{equation}
[\hat{H}_{U3}, \hat{K}_i] = -i\hbar \, \hat{p}_i. \label{eq:HU3-Galilean}
\end{equation}

This is an algebraic constraint inherited from the Galilean group's structure --- not an assumption imported from U4.

Substituting \eqref{eq:K-form} into \eqref{eq:HU3-Galilean}, and using $[\hat{H}_{U3}, t\hat{p}_i] = 0$:
\begin{equation}
m \, [\hat{H}_{U3}, \hat{x}_i] = -i\hbar \, \hat{p}_i. \label{eq:HU3-x-comm}
\end{equation}

Decomposing $\hat{H}_{U3} = T_{U3}(\hat{p}) + V_{U3}(\hat{x})$ (a working ansatz justified independently in \S\ref{sec:F5}.6) and using $[T_{U3}(\hat{p}), \hat{x}_i] = -i\hbar \, \partial T_{U3}/\partial \hat{p}_i$:
\[
m \cdot (-i\hbar) \, \frac{\partial T_{U3}(\hat{p})}{\partial \hat{p}_i} = -i\hbar \, \hat{p}_i,
\]
so
\begin{equation}
\frac{\partial T_{U3}(\hat{p})}{\partial \hat{p}_i} = \frac{\hat{p}_i}{m}. \label{eq:TU3-deriv}
\end{equation}

This is the same differential equation U4 derived. Integrating with $T_{U3}(0) = 0$:
\begin{equation}
T_{U3}(\hat{p}) = \frac{|\hat{p}|^2}{2m}. \label{eq:TU3-result}
\end{equation}

The factor of $1/2$ emerges from the chain-rule integration Jacobian, identical to U4's \S F4 derivation.

\subsection{Independent justification of the kinetic-plus-potential decomposition}

The decomposition $\hat{H}_{U3} = T_{U3}(\hat{p}) + V_{U3}(\hat{x})$ is justified at the F5 level by primitive-level translation invariance + locality (sourced to the same primitives U4 used, but invoked independently here):

\begin{itemize}
  \item \textbf{Translation invariance} of the participation graph (Primitives~03 + 06) forbids the free part of $\hat{H}_{U3}$ from depending on $\hat{x}$. The free part is $T_{U3}(\hat{p})$.
  \item \textbf{Locality} of external potentials forbids the potential part from depending on $\hat{p}$. The potential part is $V_{U3}(\hat{x})$.
  \item \textbf{Cross terms} are forbidden in the gauge-coupling-free scope.
\end{itemize}

The argument is primitive-level, sourced independently of U4's derivation chain.

\subsection{The identification}

Combining \eqref{eq:TU3-result} with \S\ref{sec:F5}.6 and restoring $\hbar$:
\begin{equation}
\hat{H}_{U3} = \frac{\hbar^2 |\hat{p}|^2}{2m} + V_{U3}(\hat{x}). \label{eq:HU3-form}
\end{equation}
Comparing with U4:
\begin{equation}
\hat{H}_{U4} = \frac{\hbar^2 |\hat{p}|^2}{2m} + V(\hat{x}). \label{eq:HU4-form}
\end{equation}
The kinetic terms are identical. The potential terms agree as $V_{U3}(\hat{x}) = V(\hat{x})$ (both denote the same external-context-supplied scalar function of position; tautological at the F5 level).

The additive vacuum-energy zero-point gauge: Stone's theorem leaves $\hat{H}_{U3}$ unique up to an additive real constant; U4's $V(\hat{x})$ has the same additive-constant ambiguity. Both are the same gauge and coincide trivially.

Therefore:
\begin{equation}
\boxed{\hat{H}_{U3} = \hat{H}_{U4} = \frac{\hbar^2 |\hat{p}|^2}{2m} + V(\hat{x}).} \label{eq:HU3-HU4}
\end{equation}

The mass $m$ is inherited per Arc~M's ``form \textsc{forced}, value \textsc{inherited}'' methodology; $\hbar$ is inherited per the dimensional-atlas Madelung anchoring. \textbf{Strategy A closes the F5 identification. Falsifier Fal-5 dispatched.}

\subsection{Strategy B (cross-check) --- Stone-theorem uniqueness on $\hat{H}_{U4}$-generated group}

The operator $\hat{H}_{U4}$ is self-adjoint on a dense domain in $\mathcal{H}$. By the spectral theorem, $\hat{H}_{U4}$ generates a strongly continuous one-parameter unitary group:
\begin{equation}
V_t := \exp(-i\hat{H}_{U4} t/\hbar). \label{eq:Vt}
\end{equation}

The group $\{V_t\}$ acts on $\mathcal{H}$ as a unitary representation of the time-translation symmetry. By the uniqueness of unitary representations of $(\mathbb{R}_t, +)$ within the Galilean group's representation on $\mathcal{H}$ (sourced to Galilean group structure), $\{V_t\}$ coincides with $\{U_t\}$ from F1.

Stone's theorem guarantees a unique self-adjoint generator of the common group $\{V_t\} = \{U_t\}$. The unique generator is $\hat{H}_{U4}$ (by construction in \eqref{eq:Vt}) and $\hat{H}_{U3}$ (by F2). Therefore $\hat{H}_{U4} = \hat{H}_{U3}$.

\textbf{Strategy B confirms Strategy A's identification by an independent route.}

\begin{theorem}[F5 --- Compatibility with U4]
Let $\hat{H}_{U3}$ denote the unique self-adjoint generator of the time-translation group $\{U_t\}$ on the U2 Hilbert space $\mathcal{H}$, supplied by Stone's theorem applied in F1 + F2. Let $\hat{H}_{U4} = \hbar^2 |\hat{p}|^2/(2m) + V(\hat{x})$ denote the kinetic-plus-potential operator derived in U4. Then, modulo the additive vacuum-energy zero-point gauge,
\[
\hat{H}_{U3} = \hat{H}_{U4} = \frac{\hbar^2 |\hat{p}|^2}{2m} + V(\hat{x})
\]
as self-adjoint operators on $\mathcal{H}$. The identification is established by Strategy A (Galilean Lie algebra closure) and cross-checked by Strategy B (Stone-uniqueness).
\end{theorem}

\subsection{Schr\"odinger evolution forced}

Combining F1, F2, F3, F4, F5:
\begin{equation}
\boxed{i\hbar \, \partial_t P(t) = \left( \frac{\hbar^2 |\hat{p}|^2}{2m} + V(\hat{x}) \right) P(t).} \label{eq:Schroedinger}
\end{equation}

This is the Schr\"odinger equation in standard form, derived structurally with no PDE-level Schr\"odinger assumption, no U4-as-derivation-input circularity, no relativistic kinematics, and no downstream gauge-coupling structure.

\section{Falsifier Resolution Table}\label{sec:falsifier-table}

\begin{center}
\begin{tabular}{p{2.0cm}p{1.0cm}p{8.5cm}p{1.8cm}}
\hline
Falsifier & Target & Dispatch argument & Section \\
\hline
Fal-1 & F1 & $\langle U_t P \mid U_t Q \rangle = \langle P \mid Q \rangle$ via translation-invariance of U2 measure on time axis & \S\ref{sec:F1}.4 \\
Fal-2 & F2 & Stone's theorem on $\{U_t\}$ produces a unique self-adjoint generator & \S\ref{sec:F2}.2 \\
Fal-3 & F3 & Linear unitary representations induce linear evolution equations & \S\ref{sec:F3}.2 \\
Fal-4 & F4 & Operator-exponential differentiation gives first-order; non-rel scope selects abelian $\mathbb{R}_t$ over Lorentz boosts & \S\ref{sec:F4}.2 \\
Fal-5 & F5 & Strategy A (Galilean Lie algebra closure) + Strategy B (Stone-uniqueness on $\hat{H}_{U4}$-generated group) & \S\ref{sec:F5} \\
Fal-6 & F2 & Stone's theorem produces a self-adjoint generator from a strongly continuous unitary group & \S\ref{sec:F2}.3 \\
Fal-7 & F1 & Pre-derivation Primitive 13 audit confirms continuous time at Phase-1 scope; $L^2$-translation continuity inherits & \S\ref{sec:F1}.6 \\
\hline
\end{tabular}
\end{center}

\textbf{All seven falsifiers dispatched.}

\section{Conditionality Ledger}\label{sec:conditionality}

The U3 verdict is \textbf{\textsc{forced}} with the following inherited conditionalities and \textbf{zero new ones}:

\begin{enumerate}
  \item Non-relativistic single-particle scope.
  \item Gauge-coupling-free scope (inherited from U4).
  \item U2 inner-product structure (both regimes).
  \item U1 participation-measure structure.
  \item U5 momentum-operator structure.
  \item U4 Hamiltonian form (as identification target only).
  \item Inherited values: mass $m$ per Arc~M; $\hbar$ per dimensional-atlas Madelung anchoring.
  \item Primitive 13 continuous-time interpretation.
  \item Galilean-group uniqueness within non-relativistic scope.
\end{enumerate}

\textbf{U3 introduces NO new structural \textsc{candidate}s.} All nine conditionalities are inherited from prior structural commitments.

\section{Final U3 Theorem}\label{sec:final-theorem}

\begin{theorem}[U3 --- Time-Translation Generator and Schr\"odinger Evolution]
Let the ED primitive stack supply: graph homogeneity (Primitive~03), bandwidth (Primitive~04), spatial gradient (Primitive~06), tension polarity (Primitive~09), and relational timing (Primitive~13). Let the participation-measure carrier (U1, \textsc{forced}), the U2 inner product (U2-Discrete + U2-Continuum, both \textsc{forced}), the momentum operator (U5, \textsc{forced}), and the kinetic-plus-potential Hamiltonian operator (U4, \textsc{forced}-conditional) supply the Hilbert space $\mathcal{H}$ and operator infrastructure. Then:

\begin{enumerate}
  \item \textbf{Time-translation representation.} The time-translation operators $(U_t P)_K(x, s) := P_K(x, s - t)$ for $t \in \mathbb{R}$ form a strongly continuous one-parameter abelian unitary group on $\mathcal{H}$.

  \item \textbf{Unique self-adjoint generator.} By Stone's theorem, there exists a unique self-adjoint operator $\hat{H}_{U3}$ on $\mathcal{H}$ such that
  \[
  U_t = \exp\!\left(-\frac{i \hat{H}_{U3} t}{\hbar}\right).
  \]

  \item \textbf{Linear first-order-in-time evolution equation.} Differentiating yields $i\hbar \, \partial_t P(t) = \hat{H}_{U3} P(t)$, linear in $P(t)$, first-order in $\partial_t$.

  \item \textbf{Compatibility with U4.} The Galilean Lie algebra forces $\hat{H}_{U3}$ to satisfy $[\hat{H}_{U3}, \hat{K}_i] = -i\hbar\hat{p}_i$ with the boost generator $\hat{K}_i = m\hat{x}_i - t\hat{p}_i$. Integration against the U5 momentum operator yields $T(\hat{p}) = |\hat{p}|^2/(2m)$ via the chain-rule Jacobian. The kinetic-plus-potential decomposition follows from translation invariance + locality. Modulo the additive vacuum-energy zero-point gauge,
  \[
  \hat{H}_{U3} = \hat{H}_{U4} = \frac{\hbar^2 |\hat{p}|^2}{2m} + V(\hat{x}).
  \]

  \item \textbf{Schr\"odinger evolution.} Combining (1)--(4):
  \[
  i\hbar \, \partial_t P(t) = \left( \frac{\hbar^2 |\hat{p}|^2}{2m} + V(\hat{x}) \right) P(t).
  \]
\end{enumerate}

The derivation invokes no PDE-level Schr\"odinger assumption, no U4-as-derivation-input circularity, no relativistic kinematics, and no downstream gauge-coupling structure. All seven falsifiers are dispatched.
\end{theorem}

\textbf{Status:} \textsc{forced}-unconditional in both discrete and continuum regimes; gauge-invariant under the U2-Continuum conformal redundancy; \textbf{U4 promoted retroactively to \textsc{forced}-unconditional}; Phase-1 program structurally complete modulo description-level continuum gauge.

\section{Downstream Consequences}\label{sec:downstream}

\subsection{U4 retroactive promotion}

U4 was closed \textsc{forced}-conditional on U3 via Framing~1. With U3 now \textsc{forced}, \textbf{U4 promotes from \textsc{forced}-conditional to \textsc{forced}-unconditional retroactively}, resolving its sibling-\textsc{candidate} conditionality.

\subsection{Schr\"odinger evolution forced-unconditional}

The Phase-1 Step~2 chain reads, post-U3:
\begin{verbatim}
Primitives 03, 04, 06, 07, 09, 13
              |
              v
        U1 (FORCED) -> P_K = sqrt(b) * exp(i pi)
              |
              v
        U2 (FORCED) -> <P|Q> inner product
              |
              v
        U5 (FORCED) -> spatial translation -> p_hat = -i hbar nabla
              |
              v
        U4 (FORCED-UNCONDITIONAL post-U3)
              |
              v
        U3 (FORCED, this paper)
              |
              v
       SCHRÖDINGER EVOLUTION (FORCED-unconditional)
\end{verbatim}

\subsection{Updated structural-foundations ledger}

\begin{center}
\begin{tabular}{cll}
\hline
\# & Theorem & Status (post-U3) \\
\hline
10 & Born rule (Gleason--Busch) & \textsc{forced}-unconditional \\
11 & U2-Discrete & \textsc{forced}-unconditional \\
12 & U2-Continuum (with conformal gauge) & \textsc{forced}-unconditional \\
13 & U5 (translation generator + momentum operator) & \textsc{forced}-unconditional \\
14 & U1 (participation-measure construction) & \textsc{forced}-unconditional \\
15 & U4 (specific Hamiltonian form) & \textbf{\textsc{forced}-unconditional} (retroactive) \\
\textbf{16} & \textbf{U3 (this paper)} & \textbf{\textsc{forced}-unconditional} \\
\hline
\end{tabular}
\end{center}

\textbf{Forced-theorem count: 15 $\to$ 16, with U4 retroactively promoted.}

\subsection{Active CANDIDATE inventory}

Pre-U3: $\{U3\}$ + continuum gauge convention. Post-U3: $\{ \, \}$ + continuum gauge convention. \textbf{The active upstream-\textsc{candidate} inventory is empty.}

\subsection{All four foundational QM postulates structurally derived}

\begin{center}
\begin{tabular}{lll}
\hline
Postulate & Pre-U3 & Post-U3 \\
\hline
Born rule & \textsc{forced}-unconditional & (unchanged) \\
Bell--Tsirelson bound & \textsc{forced}-unconditional & (unchanged) \\
Heisenberg uncertainty & \textsc{forced}-unconditional & (unchanged) \\
\textbf{Schr\"odinger evolution} & \textsc{forced}-conditional on U3 & \textbf{\textsc{forced}-unconditional} \\
\hline
\end{tabular}
\end{center}

\textbf{All four foundational postulates of non-relativistic single-particle quantum mechanics are now \textsc{forced}-unconditional structural consequences of the ED primitive ontology.}

\subsection{Phase-1 program structurally complete}

The QM-emergence Phase-1 structural-foundations program is \textbf{structurally complete} modulo the description-level continuum gauge convention. The seven-arc cycle (born\_gleason $\to$ U2 $\to$ U2-continuum $\to$ U5 $\to$ U1 $\to$ U4 $\to$ U3) closes the program.

\section{Discussion and Outlook}\label{sec:discussion}

\subsection{The kinematics-dynamics seam closed}

The structural-foundations cycle decomposes into three layers:

\begin{itemize}
  \item \textbf{Layer 1 (carrier):} U1.
  \item \textbf{Layer 2 (Hilbert-space arena):} U2-Discrete + U2-Continuum + Theorem~10 + U5.
  \item \textbf{Layer 3 (dynamical):} U4 + U3.
\end{itemize}

U5 supplies the spatial-translation generator (kinematic backbone); U3 supplies the time-translation generator (dynamical backbone); U4 supplies the form of the time-translation generator. With U3 closed, the seam between kinematics and dynamics is structurally closed: the same Stone's-theorem skeleton applies on both axes; the Galilean Lie algebra closes the loop, identifying the two generators with their U4-derived operator content.

The Schr\"odinger equation is, at root, \textbf{what falls out when Stone's theorem is applied to time-translation symmetry on the participation-measure Hilbert space, and the resulting generator is identified with the Galilean-Lie-algebra-forced kinetic-plus-potential operator.}

\subsection{Methodological completeness across seven arcs}

The structural-derivation methodology has been demonstrated against seven substantively different structural questions:

\begin{enumerate}
  \item \textbf{born\_gleason} --- Born rule via Gleason--Busch + non-contextuality.
  \item \textbf{U2-Discrete} --- sesquilinear inner product via primitive-level aggregation.
  \item \textbf{U2-Continuum} --- discrete-to-continuum lift with explicit conformal gauge.
  \item \textbf{U5} --- Stone's theorem on spatial translations $\to$ momentum operator.
  \item \textbf{U1} --- participation-measure construction via Frobenius + Cauchy functional equation.
  \item \textbf{U4} --- Galilean Lie algebra integration $\to$ kinetic-energy form, factor-of-2 Jacobian.
  \item \textbf{U3 (this paper)} --- Stone's theorem on time translations + Galilean closure $\to$ Hamiltonian operator + Schr\"odinger evolution.
\end{enumerate}

The pattern (decompose \textsc{candidate} $\to$ identify load-bearing items $\to$ close via primitive-level arguments $\to$ introduce zero new \textsc{candidate}s) holds across all seven arcs.

\subsection{What this changes}

Standard quantum mechanics adopts the Schr\"odinger equation as a postulate. The U3 result establishes that this postulate is structurally derivable: time-translation symmetry (Primitive~13) plus the U2 Hilbert space plus Stone's theorem produces the equation's existence + uniqueness + linearity + first-order character; the Galilean Lie algebra plus the U5 momentum operator forces the generator's specific kinetic-plus-potential form. No PDE-level Schr\"odinger assumption is required.

The same structural depth now applies to all four foundational QM postulates. The Hilbert-space postulate, the Born-rule postulate, the position-momentum-conjugacy postulate, and the Schr\"odinger-evolution postulate are no longer separate axioms floating above the framework --- they are derived structural consequences of the participation-graph ontology plus standard mathematical infrastructure.

\subsection{What this does not change}

The U3 result does not predict new experiments. The Schr\"odinger equation is the same equation it has always been; Event Density reproduces standard non-relativistic single-particle quantum mechanics exactly. What changes is the depth at which the answer to ``why does quantum mechanics look like this?'' can be given: in the standard formulation, ``because we postulate it''; in Event Density, ``because the participation-graph primitives plus standard mathematics force it.''

\subsection{Next steps}

\begin{enumerate}
  \item \textbf{Synthesis paper revision.} The QM-emergence Phase-1 synthesis paper requires comprehensive revision to reflect the post-U3 state.
  \item \textbf{Phase-3 (gravitational sector) extensions.} With Phase-1 closed, the natural follow-on is the Phase-3 gravitational-sector program.
  \item \textbf{Companion public-facing explainer.} The capstone explainer of the public-facing series, announcing Phase-1 closure to the general scientific audience.
  \item \textbf{Empirical-test program.} With the structural-derivation program closed, the empirical-test program (FRAP, optomechanics, BEC, retrodictions) becomes the load-bearing locus of the framework's continued development.
\end{enumerate}

\begin{thebibliography}{99}

\bibitem{phase1}
A.~Proxmire, \emph{Phase-1: QM Emergence Structural Completion}, \texttt{papers/QM\_Emergence\_Structural\_Completion/}, 2026.

\bibitem{borngleason}
A.~Proxmire, \emph{The Born Rule as a Forced Theorem of Event Density: A Gleason--Busch Reconstruction from First Principles}, \texttt{papers/Born\_Gleason/}, 2026.

\bibitem{u2}
A.~Proxmire, \emph{The Inner Product as Forced Structure in Event Density}, \texttt{papers/U2\_Inner\_Product/}, 2026.

\bibitem{u1}
A.~Proxmire, \emph{The U1 Theorem: Deriving the Participation Measure from Primitive Structure}, \texttt{papers/U1\_Participation\_Measure/}, 2026.

\bibitem{u4}
A.~Proxmire, \emph{U4: The Forced Structure of the Non-Relativistic Hamiltonian}, \texttt{papers/U4\_Hamiltonian\_Form/}, 2026.

\bibitem{u5}
A.~Proxmire, \emph{U5: The Forced Structure of Translation Symmetry and the Momentum Operator}, \texttt{papers/U5\_Translation\_Momentum/}, 2026.

\bibitem{arcm}
A.~Proxmire, \emph{Arc M: Chain-Mass Structure in Event Density}, \texttt{papers/Arc\_M/}, 2026.

\bibitem{u3arc}
A.~Proxmire, \emph{U3 arc memos}, \texttt{arcs/U3/}, 2026.

\bibitem{stone}
M.~H.~Stone, ``On one-parameter unitary groups in Hilbert space,'' \emph{Annals of Mathematics} \textbf{33}, 643--648 (1932).

\bibitem{bargmann}
V.~Bargmann, ``On unitary ray representations of continuous groups,'' \emph{Annals of Mathematics} \textbf{59}, 1--46 (1954).

\bibitem{reedsimon}
M.~Reed and B.~Simon, \emph{Methods of Modern Mathematical Physics, Volume I: Functional Analysis}, Academic Press, 1980.

\end{thebibliography}

\vspace{1em}
\noindent\rule{\linewidth}{0.4pt}
\vspace{0.2em}

\noindent\emph{Manuscript prepared April 2026. Source materials: \texttt{arcs/U3/} four-memo arc and closure memo, formalized into a single publishable document. The U3 arc is the seventh and final of the structural-foundations cycle; this paper closes the QM-emergence Phase-1 structural-foundations program at the publication-record level.}

\noindent\emph{Draft complete. Pausing for revision.}

\end{document}
